% =============================================================================
\section{NHS England Case Study}\label{sec:case_study}
% =============================================================================

The case study uses two public sources covering $1$~Aug~$2020$ -- $31$~Aug~$2022$ (the period of daily NHS regional publication). These are NHS~England COVID-19~Hospital~Activity~\cite{nhs_covid_activity}, which provides daily hospital and mechanical-ventilation bed occupancy at NHS regional resolution; and ONS Mid-Year Population Estimates~\cite{ons_mye_2021} on 2021 geography, used as the per-region PINN normaliser and the population-proportional baseline denominator. All experiments use the seven NHS England commissioning regions. The series is split chronologically by epidemic wave, so the most clinically distinct wave is held out, with training on $1$~Aug~$2020$ -- $31$~May~$2021$ (Alpha and early vaccination), validation on $1$~Jun~$30$ -- $2021$~Nov~$2021$ (Delta) for hyper-parameter selection and
early stopping, and testing on $1$~Dec~$2021$ -- $31$~Aug~$2022$ (Omicron and post-Omicron). Omicron's case-to-hospitalisation and
hospitalisation-to-critical-care ratios differ markedly from those of earlier waves, providing a regime-shift stress test.

For the allocation problem, baseline capacity $\Gamma_r$ is \emph{measured}
as each region's Delta-wave peak MV-bed occupancy $\times 1.05$, whereas
the per-region surge cap $K_r=0.5\,\Gamma_r$, total budget
$B=0.20\sum_r \Gamma_r\approx213$~beds (the surge headroom above the
$\sum_r\Gamma_r\approx1065$-bed measured baseline, not existing beds),
mutual-aid travel cap $\tau=240$~min, and
uncapped outbound transfers ($\kappa=\infty$) are
planning assumptions. The surge $b_r$ is the model's decision, with the
extra MV beds a region stands up above its measured baseline. The
forecaster, budget $B$, tail weight $\lambda_3$, travel cap $\tau$, and
transfer cap $\kappa$ are each swept in \Cref{ssec:e5_e6}. Forecast
metrics are macro-averaged across
the seven regions and the $32$ rolling-origin evaluations spanning the
Omicron test period; the full solver comparison is reported at the first
Omicron decision origin ($29$~Dec~$2021$), with an
all-origin check in \Cref{ssec:e5_e6}.
These are stress-test settings, not NHS standards: the $5\%$ uplift is
a headroom proxy, $K_r$ prevents unlimited concentration, and the
$20\%$ budget and $240$-min window form a generous headline case.
Deployment requires staffed-bed inventories, local expansion and
transfer limits, network travel times, and joint residuals for scenario
probabilities; our proxy supports policy comparison, not deployable
bed-count claims.
